% Vietnamese translation for OI Public Library
% Source-PDF: 英文资料 ENGLISH MATERIALS/Longest common subsequence.pdf
\documentclass[11pt]{article}
\usepackage{oipl-vn}

\newcommand{\LCS}{\operatorname{LCS}}
\newcommand{\LIS}{\operatorname{LIS}}
\newcommand{\SC}{\operatorname{SC}}

\title{Dãy Con Chung Dài Nhất}
\author{Tài liệu bài giảng}
\date{}

\begin{document}
\maketitle

\origtitle{Longest Common Subsequence}
\sourcepath{English Materials/Longest common subsequence.pdf}

\section{Định nghĩa}

Dãy con chung dài nhất, hay \(\LCS\), của hai xâu \(S_1\) và \(S_2\) là dãy
con dài nhất xuất hiện trong cả hai xâu. Ví dụ, với hai xâu được căn chỉnh:
\[
\begin{array}{r*{13}{c}}
S_1:& A&-&A&T&-&G&G&C&C&-&A&T&A\\
S_2:& A&T&A&T&A&A&T&T&C&T&A&T&-
\end{array}
\]
một \(\LCS\) là \texttt{AATCAT}, có độ dài \(6\). Lời giải không nhất thiết
duy nhất. Với cặp \((\texttt{ATTA},\texttt{ATAT})\), các lời giải
\(\texttt{ATT}\) và \(\texttt{ATA}\) đều là dãy con chung dài nhất. Nói chung,
với hai xâu bất kỳ có thể tồn tại nhiều lời giải.

\section{Liên hệ với căn chỉnh và khoảng cách chỉnh sửa}

\paragraph{Định lý.}
Nếu cho điểm \(1\) cho mỗi ký tự khớp, và điểm \(0\) cho mỗi ký tự không khớp
hoặc ô trống, thì các ký tự khớp trong một căn chỉnh có giá trị cực đại tạo
thành một \(\LCS\).

Vì cách làm này dùng quy hoạch động tổng quát, độ phức tạp là \(O(nm)\), với
\(n=|S_1|\) và \(m=|S_2|\). Bài toán xâu con chung dài nhất theo nghĩa đoạn
liên tiếp, ngược lại, có lời giải \(O(n+m)\). Dãy con phức tạp hơn nhiều so
với xâu con liên tiếp.

Trong ví dụ ở trên, căn chỉnh tối ưu có ba phép chèn, một phép xóa và ba phép
thay thế, nên khoảng cách chỉnh sửa là \(7\) nếu thay thế có chi phí \(1\).
Nhưng mỗi phép thay thế có thể được thực hiện bằng một phép xóa rồi một phép
chèn. Nếu chi phí thay thế là \(2\), còn chèn và xóa đều có chi phí \(1\), thì
khoảng cách chỉnh sửa nhỏ nhất \(D\) liên hệ với độ dài \(\LCS\), ký hiệu là
\(L\), bởi:
\[
  D=m+n-2L.
\]
Với ví dụ trên, \(n=10\), \(m=12\), \(L=6\), nên \(D=10\): gồm \(6\) phép chèn
và \(4\) phép xóa.

\section{Tính trực tiếp bằng quy hoạch động}

Ta có thể tính \(\LCS\) trực tiếp bằng quy hoạch động. Cách này hiệu quả hơn
về hằng số, dù độ phức tạp tiệm cận vẫn là \(O(nm)\). Gọi \(L(i,j)\) là độ
dài \(\LCS\) của hai tiền tố \(S_1[1..i]\) và \(S_2[1..j]\). Khi đó:
\[
\begin{aligned}
L(0,0)&=0,\\
L(i,0)&=0,\\
L(0,j)&=0,
\end{aligned}
\]
và với \(i,j>0\):
\[
L(i,j)=
\begin{cases}
1+L(i-1,j-1), & S_1(i)=S_2(j),\\
\max\{L(i,j-1),L(i-1,j)\}, & S_1(i)\ne S_2(j).
\end{cases}
\]
Nếu lưu thêm con trỏ truy vết trong ma trận, ta có thể khôi phục một hoặc
nhiều dãy con chung dài nhất.

\paragraph{Ví dụ.}
Với \(S_1=\texttt{ATCAAT}\) và \(S_2=\texttt{TAATCATA}\), một giá trị
\(\LCS\) là \(5\), nên khoảng cách chỉnh sửa với chi phí thay thế bằng \(2\)
là:
\[
  6+8-2\cdot5=4.
\]
Bảng quy hoạch động trong tài liệu gốc ghi các giá trị \(L(i,j)\) cho từng
tiền tố; hàng cuối cho thấy độ dài cực đại là \(5\).

\section{Một hướng nhanh hơn thông qua LIS}

Ta xét một bài toán phụ. Cho \(\pi\) là một dãy gồm \(n\) số nguyên, không
nhất thiết phân biệt.

\paragraph{Dãy con tăng.}
Dãy con tăng (\(\LIS\) theo nghĩa \textit{increasing subsequence}) của
\(\pi\) là một dãy con có giá trị tăng nghiêm ngặt từ trái sang phải. Ví dụ:
\[
  \pi=(5,3,4,4,9,6,2,1,8,7,10)
\]
có các dãy con tăng như \((3,4,6,8,10)\) và \((5,9,10)\). Dãy con tăng dài
nhất là dãy con tăng có độ dài cực đại.

\paragraph{Dãy con giảm và phủ.}
Dãy con giảm ở đây là dãy con không tăng. Ví dụ:
\[
  (5,4,4,2,1).
\]
Một \term{phủ} là một tập các dãy con giảm rời nhau của \(\pi\) sao cho mọi
phần tử của \(\pi\) đều nằm trong một dãy. Kích thước phủ \(c\) là số dãy con
giảm trong phủ. Ví dụ:
\[
\pi=(5,3,4,9,6,2,1,8,7)
\]
có phủ
\[
\{(5,3,2,1),(4),(9,6),(8,7)\},
\]
nên \(c=4\). Phủ nhỏ nhất là phủ có \(c\) nhỏ nhất.

\paragraph{Bổ đề.}
Nếu \(I\) là một dãy con tăng của \(\pi\) có độ dài bằng kích thước của một
phủ \(C\) của \(\pi\), thì \(I\) là dãy con tăng dài nhất, và \(C\) là phủ nhỏ
nhất.

Thật vậy, một dãy con tăng không thể chứa nhiều hơn một phần tử từ cùng một
dãy con giảm. Nếu
\[
  C=C_1\cup C_2\cup\cdots\cup C_k,
\]
thì mọi dãy con tăng có độ dài nhiều nhất \(k=|C|\). Ngược lại, nếu có một
phủ nhỏ hơn \(C\), dãy con tăng \(I\) sẽ buộc phải lấy hơn một phần tử từ một
dãy con giảm nào đó của phủ ấy, điều không thể xảy ra.

\section{Xây dựng phủ trong thời gian log-tuyến tính}

Thuật toán tham lam xây dựng phủ như sau. Duyệt \(\pi\) từ trái sang phải. Với
mỗi số hiện tại, chèn nó vào dãy con giảm nằm trái nhất trong các dãy đã tạo
sao cho tính không tăng vẫn được giữ. Nếu không chèn được vào dãy nào, bắt đầu
một dãy con giảm mới.

Để đạt \(O(n\log n)\), lưu một danh sách \(L\) gồm phần tử cuối của mỗi dãy
con giảm đang xây dựng, kèm định danh của dãy. Danh sách này luôn được sắp
theo giá trị tăng dần. Với phần tử mới \(x_i\), tìm bằng nhị phân phần tử đầu
tiên trong \(L\) có giá trị lớn hơn \(x_i\). Nếu tìm thấy, nối \(x_i\) vào
dãy tương ứng và thay giá trị cuối trong \(L\) bằng \(x_i\). Nếu không tìm
thấy, tạo một dãy mới và thêm \((x_i,j)\) vào \(L\).

\paragraph{Định lý.}
Phủ tham lam có thể xây dựng trong \(O(n\log n)\). Đồng thời, một dãy con tăng
dài nhất và một phủ nhỏ nhất cũng có thể được xây dựng trong \(O(n\log n)\).

Với
\[
  \pi=(5,3,4,9,6,2,1,8,7,10),
\]
thuật toán tạo các dãy:
\[
  D_1=(5,3,2,1),\quad D_2=(4),\quad D_3=(9,6),\quad
  D_4=(8,7),\quad D_5=(10).
\]
Phần giá trị của danh sách \(L\) lần lượt đi qua:
\[
\begin{gathered}
(5),(3),(3,4),(3,4,9),(3,4,6),(2,4,6),\\
(1,4,6),(1,4,6,8),(1,4,6,7),(1,4,6,7,10).
\end{gathered}
\]

\section{Quy LCS về LIS}

Cho hai dãy \(S_1\) và \(S_2\). Gọi \(r_i\) là số lần xuất hiện của ký tự thứ
\(i\) trong \(S_1\) bên trong \(S_2\). Ví dụ:
\[
  S_1=\texttt{abacx},\qquad S_2=\texttt{baabca}.
\]
Khi đó:
\[
  r(1)=3,\quad r(2)=2,\quad r(3)=3,\quad r(4)=1,\quad r(5)=0.
\]

Với mỗi ký tự phân biệt \(x\) trong \(S_1\), định nghĩa \(\operatorname{list}(x)\)
là danh sách các vị trí của \(x\) trong \(S_2\), viết theo thứ tự giảm dần. Ví
dụ:
\[
  \operatorname{list}(\texttt{a})=(6,3,2),\quad
  \operatorname{list}(\texttt{b})=(4,1),\quad
  \operatorname{list}(\texttt{c})=(5),\quad
  \operatorname{list}(\texttt{x})=\varnothing.
\]

Định nghĩa \(\Pi(S_1,S_2)\) là dãy thu được bằng cách nối
\(\operatorname{list}(s_i)\) với \(i=1,2,\ldots,n\), trong đó \(n=|S_1|\) và
\(s_i\) là ký tự thứ \(i\) của \(S_1\). Trong ví dụ:
\[
  \Pi(S_1,S_2)=(6,3,2,4,1,6,3,2,5).
\]

\paragraph{Định lý.}
Mọi dãy con tăng \(I\) của \(\Pi(S_1,S_2)\) xác định một dãy con chung của
\(S_1\) và \(S_2\) có cùng độ dài, và ngược lại. Vì vậy một \(\LCS\) của
\(S_1\) và \(S_2\) tương ứng với một \(\LIS\) của \(\Pi(S_1,S_2)\).

Trong ví dụ trên, các dãy con tăng dài nhất của
\(\Pi(S_1,S_2)\) có thể dùng làm chỉ số truy cập vào \(S_2\) để thu được các
\(\LCS\):
\[
  (1,2,5)\mapsto\texttt{bac},\qquad
  (2,3,5)\mapsto\texttt{aac},\qquad
  (3,4,6)\mapsto\texttt{aba}.
\]

\end{document}
